% =====================================================
% CHAPITRE 3 : COMPRÉHENSION DES DONNÉES (CRISP-DM Phase 2)
% =====================================================
\chapter{Compréhension des données}

% ---------------------------------------------------
\section*{Introduction}
\addcontentsline{toc}{section}{Introduction}
% ---------------------------------------------------

La phase de compréhension des données constitue le socle de tout projet de data science.
Elle précède et conditionne chaque décision de modélisation : un modèle ne peut être meilleur
que les données sur lesquelles il est entraîné. Avant d'appliquer le moindre algorithme,
il est indispensable de saisir la structure, les limites et les opportunités du corpus
disponible. Ce chapitre correspond au notebook \texttt{01\_data\_understanding.ipynb} du
projet et présente le jeu de données, son analyse exploratoire, l'étude de l'existant,
ainsi que l'état de l'art des approches de détection d'anomalies financières.

% ---------------------------------------------------
\section{Présentation du jeu de données}
% ---------------------------------------------------

\subsection{Origine et description}

Le jeu de données utilisé dans ce projet provient d'une \textbf{mission d'audit IT réalisée
au sein de PwC Tunisie}. Il regroupe des transactions financières issues du système
d'information d'un client. Pour des raisons de confidentialité, toutes les informations
permettant d'identifier les parties concernées ont été anonymisées.

\subsection{Caractéristiques générales}

Le corpus est constitué d'environ \textbf{200~000 transactions financières} enregistrées
sur une période couvrant 743 heures (soit environ 31 jours). Parmi ces transactions, seules
\textbf{258 sont frauduleuses}, soit un taux de fraude de \textbf{0,129\,\%}.
Ce déséquilibre extrême (ratio 1:774) constitue le défi central du projet.

\subsection{Description des variables}

\begin{table}[H]
\centering
\caption{Description des variables du jeu de données}
\label{tab:variables}
\begin{tabular}{lll}
\toprule
\textbf{Variable} & \textbf{Type} & \textbf{Description} \\
\midrule
\texttt{step}           & Entier    & Pas de temps (1 = 1 heure, max 743) \\
\texttt{type}           & Catégorie & Type d'opération (5 modalités) \\
\texttt{amount}         & Réel      & Montant de la transaction \\
\texttt{nameOrig}       & Chaîne    & Identifiant du compte émetteur \\
\texttt{oldbalanceOrg}  & Réel      & Solde émetteur avant transaction \\
\texttt{newbalanceOrig} & Réel      & Solde émetteur après transaction \\
\texttt{nameDest}       & Chaîne    & Identifiant du compte destinataire \\
\texttt{oldbalanceDest} & Réel      & Solde destinataire avant transaction \\
\texttt{newbalanceDest} & Réel      & Solde destinataire après transaction \\
\texttt{isFraud}        & Binaire   & \textbf{Variable cible} : 1 = fraude \\
\texttt{isFlaggedFraud} & Binaire   & Indicateur système naïf existant \\
\bottomrule
\end{tabular}
\end{table}

\subsection{Types de transactions}

Cinq types de transactions sont présents dans le jeu de données :

\begin{table}[H]
\centering
\caption{Distribution des types de transactions et taux de fraude}
\label{tab:types_transactions}
\begin{tabular}{lrr}
\toprule
\textbf{Type} & \textbf{Fréquence} & \textbf{Taux de fraude} \\
\midrule
CASH\_IN   & $\sim$35\,\% & 0,000\,\% \\
CASH\_OUT  & $\sim$35\,\% & 0,181\,\% \\
DEBIT      & $\sim$1\,\%  & 0,000\,\% \\
PAYMENT    & $\sim$22\,\% & 0,000\,\% \\
TRANSFER   & $\sim$7\,\%  & 0,775\,\% \\
\bottomrule
\end{tabular}
\end{table}

\textbf{Observation clé :} \textit{seuls les types TRANSFER et CASH\_OUT présentent des cas
de fraude}. Cette information sera directement exploitée lors de l'ingénierie de variables
(création de la feature \texttt{is\_transfer\_or\_cashout}).

% ---------------------------------------------------
\section{Analyse exploratoire des données (EDA)}
% ---------------------------------------------------

\subsection{Déséquilibre des classes}

Le déséquilibre est extrême : les transactions frauduleuses représentent \textbf{0,129\,\%}
du corpus. Un modèle prédisant systématiquement « légitime » atteindrait une exactitude de
99,87\,\% tout en ne détectant \textbf{aucune fraude}. C'est pourquoi l'exactitude
(\textit{accuracy}) est une métrique totalement inadaptée à ce problème, et que les
métriques de rappel, précision et F1-score ont été retenues comme indicateurs de succès.

\subsection{Distribution des montants}

La variable \texttt{amount} présente une asymétrie (\textit{skewness}) très élevée
(\textbf{30,80}), indiquant une distribution fortement concentrée sur de petites valeurs
avec de rares transactions à montant très élevé.

\begin{table}[H]
\centering
\caption{Statistiques descriptives de la variable \texttt{amount}}
\label{tab:amount_stats}
\begin{tabular}{ll}
\toprule
\textbf{Statistique} & \textbf{Valeur} \\
\midrule
Minimum     & 0,00 \\
Médiane     & $\sim$ 74~000 \\
Moyenne     & $\sim$ 179~862 \\
Maximum     & $\sim$ 92~445~517 \\
Skewness    & 30,80 \\
\bottomrule
\end{tabular}
\end{table}

Cette propriété justifie l'application d'une transformation logarithmique
$\log(1 + x)$ lors de la préparation des données pour corriger la skewness et
améliorer les performances des modèles linéaires.

\subsection{Analyse temporelle}

L'analyse de la variable \texttt{step} révèle une \textbf{saisonnalité des fraudes}.
Les heures à risque élevé identifiées dans l'EDA sont :
$\{0, 1, 2, 3, 4, 5, 6, 7, 8, 9, 23\}$ --- soit les heures nocturnes et
de début de matinée. Le ratio de fraude durant ces heures est de \textbf{10,63 fois
supérieur} à celui des heures ordinaires. Cette observation conduit à la création d'une
variable indicatrice \texttt{high\_risk\_hour}.

\subsection{Analyse des soldes et détection de leakage}

L'analyse des variables de soldes révèle des incohérences pour les transactions frauduleuses :
\texttt{oldbalanceOrg}, \texttt{newbalanceOrig}, \texttt{oldbalanceDest} et
\texttt{newbalanceDest} présentent des soldes nuls anormaux. Ces variables sont des
\textbf{indicateurs trop directs de la fraude} et pourraient causer une fuite
d'information (\textit{data leakage}) si elles étaient directement incluses dans les
modèles. À la place, une variable dérivée \texttt{balance\_diff\_orig} est créée, avec
une corrélation de \textbf{0,3662} avec la variable cible \texttt{isFraud}.

\subsection{Baseline du système existant}

Le système existant (\texttt{isFlaggedFraud}) sert de référence basse. Ses performances
illustrent l'inadéquation des approches par règles simples :

\begin{table}[H]
\centering
\caption{Performances de la baseline \texttt{isFlaggedFraud}}
\label{tab:baseline_perf}
\begin{tabular}{lr}
\toprule
\textbf{Métrique} & \textbf{Valeur} \\
\midrule
Recall    & 0,0039 \\
Precision & 1,0000 \\
F1-Score  & 0,0077 \\
\bottomrule
\end{tabular}
\end{table}

\textit{Ce système ne détecte qu'une infime fraction des fraudes réelles --- une seule
transaction sur 258 ---, confirmant la nécessité d'une approche par apprentissage
automatique.}

% ---------------------------------------------------
\section{Étude de l'existant et état de l'art}
% ---------------------------------------------------

\subsection{Approches supervisées}

Les approches supervisées nécessitent des étiquettes (\textit{labels}) pour chaque
transaction. Parmi les plus utilisées pour la détection de fraude :

\begin{itemize}
  \item \textbf{Régression logistique} : baseline linéaire, faible coût computationnel,
    bonne interprétabilité.
  \item \textbf{Forêt aléatoire} : ensembliste, robuste au bruit, gère bien les
    interactions non linéaires entre variables.
  \item \textbf{Gradient Boosting (XGBoost, LightGBM)} : performances état de l'art
    sur données tabulaires.
  \item \textbf{Réseaux de neurones feedforward} : puissants mais nécessitent beaucoup
    de données labellisées.
\end{itemize}

\subsection{Approches non supervisées}

Lorsque les étiquettes sont rares ou indisponibles, les approches non supervisées
permettent de détecter des patterns inhabituels sans connaissance préalable des fraudes :

\begin{itemize}
  \item \textbf{Isolation Forest} : isole les anomalies par partitions aléatoires
    successives.
  \item \textbf{One-Class SVM} : délimite la frontière de la normalité dans l'espace
    des features.
  \item \textbf{Autoencoder} : apprend à reconstruire les données normales ; une erreur
    de reconstruction élevée signale une anomalie.
  \item \textbf{Autoencoders variationnels (VAE)} : modélisation probabiliste de la
    distribution normale, plus robuste pour les données continues.
\end{itemize}

\subsection{Approches retenues pour ce projet}

Pour ce projet, trois approches complémentaires ont été retenues, couvrant le spectre
supervisé / non supervisé :

\begin{enumerate}
  \item \textbf{Régression logistique} --- baseline supervisée simple, pour établir
    une référence linéaire solide.
  \item \textbf{Forêt aléatoire} --- référence supervisée robuste, capable de capturer
    les interactions non linéaires entre variables.
  \item \textbf{Autoencoder} --- approche non supervisée orientée reconstruction, utile
    lorsque les labels sont limités ou lorsqu'on souhaite modéliser les patterns normaux
    sans dépendre des exemples de fraude.
\end{enumerate}

Ce choix vise à couvrir plusieurs angles d'attaque du problème, afin de disposer d'une
comparaison objective et d'une complémentarité méthodologique.

% ---------------------------------------------------
\section*{Conclusion}
\addcontentsline{toc}{section}{Conclusion}
% ---------------------------------------------------

L'analyse exploratoire a mis en évidence les caractéristiques clés du jeu de données :
déséquilibre extrême des classes (1:774), asymétrie des montants (skewness = 30,80),
saisonnalité temporelle des fraudes (ratio 10,63× la nuit) et risque de fuite d'information
sur les variables de soldes. Ces observations ne sont pas de simples curiosités statistiques
--- elles orientent directement chaque décision de préparation et de modélisation décrite
dans les chapitres suivants. L'inadéquation flagrante du système de règles existant
(F1 = 0,0077) confirme l'utilité d'une approche par machine learning. Le chapitre suivant
décrit en détail les opérations de préparation des données.
